\documentclass[11pt,a4paper]{article}

\usepackage[utf8]{inputenc}
\usepackage{amsmath,amssymb,amsfonts}
\usepackage{booktabs,tabularx,longtable}
\usepackage{enumitem}
\usepackage{geometry}
\usepackage{hyperref}
\usepackage{xcolor}
\usepackage{listings}
\usepackage{fancyhdr}
\usepackage{titlesec}

\geometry{margin=2.2cm,top=2.5cm,bottom=2.5cm}

\hypersetup{
  colorlinks=true,
  linkcolor=blue!60!black,
  citecolor=blue!60!black,
  urlcolor=blue!60!black,
  pdftitle={ISLS v2.2--v3.0 Unified Specification},
  pdfauthor={Sebastian Klemm},
}

\lstset{
  basicstyle=\ttfamily\scriptsize,
  breaklines=true,
  frame=single,
  backgroundcolor=\color{gray!8},
  keywordstyle=\color{blue!70!black}\bfseries,
  stringstyle=\color{red!60!black},
  commentstyle=\color{green!50!black}\itshape,
  numbers=left,
  numberstyle=\tiny\color{gray},
  numbersep=5pt,
  tabsize=4,
  showstringspaces=false,
}

\titleformat{\section}{\Large\bfseries}{\thesection}{1em}{}
\titleformat{\subsection}{\large\bfseries}{\thesubsection}{1em}{}
\titleformat{\subsubsection}{\normalsize\bfseries}{\thesubsubsection}{1em}{}

\pagestyle{fancy}
\fancyhf{}
\fancyhead[L]{\small ISLS v2.2--v3.0 Unified Spec}
\fancyhead[R]{\small Klemm, 2026}
\fancyfoot[C]{\thepage}

\begin{document}

\begin{titlepage}
\begin{center}
\vspace*{1cm}
{\Huge\bfseries ISLS Unified Specification}\\[0.4cm]
{\LARGE\bfseries v2.2: Type-Aware LLM + Multi-Domain}\\[0.2cm]
{\LARGE\bfseries v3.0: Self-Defining Norm System}\\[0.2cm]
{\LARGE\bfseries + Chat Agent + Studio Dashboard}\\[1cm]
{\Large Master Implementation Document}\\[1.5cm]
{\large Sebastian Klemm}\\[0.3cm]
{\small March 2026}\\[1.5cm]

\begin{tabular}{ll}
\textbf{Repository:} & \texttt{graphity/isls} \\
\textbf{Builds on:} & v2.1 (Hypercube Decomposer + Render Loop) \\
\textbf{New crates:} & \texttt{isls-norms}, \texttt{isls-chat} \\
\textbf{Modified:} & \texttt{isls-renderloop}, \texttt{isls-orchestrator}, \\
 & \texttt{isls-hypercube}, \texttt{isls-studio}, \\
 & \texttt{isls-agent}, \texttt{isls-gateway}, \texttt{isls-cli} \\
\textbf{License:} & MIT \\
\end{tabular}

\vfill
{\small
This document combines all remaining work into one specification:\\[0.2cm]
\textbf{Part I}: Fix LLM enrichment (type context), add 2 more domains,\\
harden output quality. The output must be \emph{good}, not just compilable.\\[0.2cm]
\textbf{Part II}: Self-defining norm system where software building blocks\\
are discovered automatically from crystal accumulation.\\[0.2cm]
\textbf{Part III}: Chat agent that translates natural language into norm\\
compositions. Homer says what he wants. The system builds it.\\[0.2cm]
\textbf{Part IV}: Studio dashboard with chat, project management,\\
norm explorer, file browser, and convergence visualization.\\[0.4cm]
Implementation order: Part I $\to$ Part II $\to$ Part III $\to$ Part IV.\\
Each part must compile and pass all tests before the next begins.
}
\end{center}
\end{titlepage}

\tableofcontents
\newpage

% %%%%%%%%%%%%%%%%%%%%%%%%%%%%%%%%%%%%%%%%%%%%%%%%%%%%%%%%%%%%
\part{Type-Aware LLM Enrichment + Multi-Domain}
\label{part:v22}
% %%%%%%%%%%%%%%%%%%%%%%%%%%%%%%%%%%%%%%%%%%%%%%%%%%%%%%%%%%%%

\section{Type Context Collector}

\subsection{Problem}

When \texttt{should\_enrich} was enabled in v2.1, the LLM wrote
code referencing non-existent fields (\texttt{payload.price} instead
of \texttt{payload.unit\_price\_cents}), wrong error constructors
(\texttt{AppError::ValidationError("str".into())} but it takes
\texttt{Vec<String>}), and wrong imports (\texttt{use log::} instead
of \texttt{tracing}).

\subsection{Solution: TypeContext}

Create \texttt{isls-renderloop/src/type\_context.rs}:

\begin{lstlisting}[language=C]
/// Extracts exact type definitions from generated files
/// so the LLM sees real field names, types, and signatures.
pub struct TypeContext {
    pub models: Vec<ModelDef>,
    pub error_enum: String,
    pub pagination_types: String,
    pub auth_types: String,
    pub query_signatures: HashMap<String, Vec<String>>,
}

pub struct ModelDef {
    pub entity_name: String,
    pub main_struct: String,
    pub create_payload: String,
    pub update_payload: String,
    pub validation_impl: String,
}

impl TypeContext {
    /// Build from generated output directory.
    /// Call AFTER Pass 0, BEFORE Pass 1.
    pub fn from_output_dir(output_dir: &Path) -> Result<Self> {
        // Read backend/src/errors.rs -> extract AppError enum
        // Read backend/src/pagination.rs -> extract PaginationParams, PaginatedResponse
        // Read backend/src/auth.rs -> extract AuthUser, Claims
        // Read backend/src/models/*.rs -> extract all structs per entity
        // Read backend/src/database/*_queries.rs -> extract fn signatures
    }

    /// Build prompt context for a specific entity.
    pub fn prompt_for_entity(&self, entity_name: &str) -> String {
        // Returns formatted string with:
        // ## EXACT TYPE DEFINITIONS
        // ### Error Type   (AppError enum)
        // ### Pagination   (PaginationParams, PaginatedResponse)
        // ### Auth          (AuthUser)
        // ### Model         (Entity struct, CreatePayload, UpdatePayload)
        // ### Query Functions (available fn signatures)
        // ### RULES
        // - Use tracing, not log
        // - ValidationError takes Vec<String>
        // - PaginationParams fields are Option<i64>
        // - Use EXACT field names
    }
}
\end{lstlisting}

\subsection{Helper Functions}

\begin{lstlisting}[language=C]
/// Extract a struct definition (derives + fields) by name.
fn extract_struct_full(content: &str, name: &str) -> String;
/// Extract pub [async] fn signatures (no body).
fn extract_fn_signatures(content: &str) -> Vec<String>;
/// Extract an enum definition by name.
fn extract_enum_definition(content: &str, name: &str) -> String;
/// Extract field names from a struct definition string.
fn extract_field_names(struct_def: &str) -> Vec<String>;
/// Extract field.access patterns from code (for hallucination check).
fn extract_field_accesses(code: &str) -> Vec<String>;
\end{lstlisting}

All use simple string matching (find struct name, count braces).
No Rust parser dependency.

\subsection{Function-Level Replacement}

The LLM enriches ONE function at a time. The file is NOT replaced.

\begin{lstlisting}[language=C]
/// Replace a single function in a file.
/// Finds fn by name, counts braces to find body, replaces body only.
fn replace_function_in_file(
    file_content: &str,
    function_name: &str,
    new_function: &str,
) -> String {
    // 1. Find "pub async fn {name}(" or "pub fn {name}("
    // 2. Walk back to doc comments
    // 3. Walk forward, count { } to find matching close
    // 4. Replace that range with new_function
    // 5. Keep everything else unchanged
}
\end{lstlisting}

\subsection{Enrichment Validation}

Before inserting LLM output:

\begin{lstlisting}[language=C]
fn validate_enriched_function(
    code: &str,
    type_ctx: &TypeContext,
    entity_name: &str,
) -> bool {
    // 1. Not empty
    // 2. No markdown fences (```)
    // 3. Balanced braces
    // 4. No `use log::`
    // 5. Contains `fn `
    // 6. No hallucinated field names (check against model struct fields)
    //    - Extract .field accesses from code
    //    - Compare against known fields from TypeContext
    //    - Skip common non-entity fields (pool, params, user, id, etc.)
    //    - If unknown field found -> REJECT, keep template version
}
\end{lstlisting}

\subsection{Re-enable should\_enrich (Selective)}

\begin{lstlisting}[language=C]
fn should_enrich(path: &str) -> bool {
    if path.contains("services/") && !path.ends_with("mod.rs") { return true; }
    if path.contains("auth_routes") { return true; }
    false // everything else: protected
}
\end{lstlisting}

\subsection{Modified DomainLogic Pass}

\begin{lstlisting}[language=C]
fn pass_domain_logic(&mut self, artifacts: &mut Vec<Artifact>, spec: &AppSpec) -> Result<()> {
    let type_ctx = TypeContext::from_output_dir(&self.output_dir)?;

    for artifact in artifacts.iter_mut() {
        if !should_enrich(artifact.path()) { continue; }
        let entity = extract_entity_from_path(artifact.path())?;
        let current = artifact.content();
        let functions = find_enrichable_functions(current);

        for func in &functions {
            let prompt = format!(
                "Fill in the body of this Rust service function.\n\n\
                 {}\n\n\
                 ## Function (keep EXACT signature)\n```rust\n{}\n```\n\n\
                 ## Business Rule\n{}\n\n\
                 Return the COMPLETE function (signature + body). No use statements.",
                type_ctx.prompt_for_entity(&entity),
                func.signature,
                func.business_rule.as_deref().unwrap_or("Standard CRUD"),
            );
            let response = self.oracle.call(&prompt, 2048)?;
            let enriched = extract_function_from_response(&response, &func.name);
            if validate_enriched_function(&enriched, &type_ctx, &entity) {
                let new_content = replace_function_in_file(current, &func.name, &enriched);
                artifact.set_content(new_content);
            }
        }
    }
    Ok(())
}
\end{lstlisting}

\subsection{TypeContext Tests}

10 tests:
\begin{itemize}
\item Extract model struct from sample file
\item Extract fn signatures from sample queries file
\item Extract enum definition (AppError)
\item \texttt{prompt\_for\_entity} returns non-empty with all sections
\item \texttt{validate\_enriched\_function} accepts good code
\item \texttt{validate\_enriched\_function} rejects hallucinated fields
\item \texttt{validate\_enriched\_function} rejects markdown fences
\item \texttt{validate\_enriched\_function} rejects \texttt{use log::}
\item \texttt{replace\_function\_in\_file} replaces target only
\item \texttt{replace\_function\_in\_file} keeps other functions intact
\end{itemize}

% ----------------------------------------------------------
\section{Multi-Domain: E-Commerce + Project Management}

\subsection{ecommerce.toml}

\begin{lstlisting}
[app]
name = "ecommerce-platform"
description = "Online store with catalog, cart, and order management"
[app.modules]
catalog = "Product listings, categories, search, reviews"
cart = "Shopping cart, add/remove items, quantity updates"
orders = "Checkout, order processing, payment tracking, shipping"
customers = "Customer accounts, addresses, order history"
auth = "JWT authentication, customer registration, admin roles"
[backend]
language = "rust"
framework = "actix-web"
database = "postgresql"
auth_method = "jwt"
[frontend]
type = "spa"
framework = "vanilla"
styling = "minimal"
[deployment]
containerized = true
compose = true
[constraints]
max_crates = 1
test_coverage = "integration"
evidence_chain = true
\end{lstlisting}

\subsection{project.toml}

\begin{lstlisting}
[app]
name = "project-tracker"
description = "Project management with tasks, sprints, and teams"
[app.modules]
projects = "Project creation, settings, team assignment, status"
tasks = "Task creation, assignment, priority, status workflow"
sprints = "Sprint planning, backlog grooming, velocity tracking"
teams = "Team members, roles, workload view"
auth = "JWT authentication, role-based access (admin, lead, member)"
[backend]
language = "rust"
framework = "actix-web"
database = "postgresql"
auth_method = "jwt"
[frontend]
type = "spa"
framework = "vanilla"
styling = "minimal"
[deployment]
containerized = true
compose = true
[constraints]
max_crates = 1
test_coverage = "integration"
evidence_chain = true
\end{lstlisting}

\subsection{Domain Templates in isls-hypercube}

Add \texttt{ecommerce\_domain()} and \texttt{project\_management\_domain()}
to the DomainRegistry alongside \texttt{warehouse\_domain()}.

\textbf{E-Commerce entities:} Product (name, slug, price\_cents, compare\_price,
sku, inventory\_count, category\_id, is\_published, image\_url, weight\_grams),
Category (name, slug, parent\_id, position, is\_active),
Customer (email, password\_hash, first\_name, last\_name, phone, is\_active),
Address (customer\_id, label, street, city, state, postal\_code, country, is\_default),
Cart (customer\_id, session\_id, status), CartItem (cart\_id, product\_id, quantity,
unit\_price\_cents), Order (order\_number, customer\_id, status, subtotal, tax,
shipping, total, payment\_status, shipping\_address\_id), OrderLine (order\_id,
product\_id, quantity, unit\_price\_cents), Review (product\_id, customer\_id,
rating 1--5, title, body, is\_approved).

Business rules: cart-to-order conversion, inventory check, order state machine,
review moderation, duplicate review prevention.

\textbf{Project Management entities:} Project (name, description, owner\_id,
status, start\_date, target\_end\_date), Sprint (project\_id, name, goal,
start\_date, end\_date, status), Task (project\_id, sprint\_id, title,
description, assignee\_id, reporter\_id, status, priority, estimate\_hours,
actual\_hours, due\_date, position), Comment (task\_id, author\_id, body),
Label (project\_id, name, color), TaskLabel (task\_id, label\_id),
TeamMember (user\_id, project\_id, role), User (email, password\_hash,
display\_name, avatar\_url, is\_active).

Business rules: task state machine (todo$\to$in\_progress$\to$in\_review$\to$done,
any$\to$blocked), sprint velocity calculation, role-based task assignment.

% ----------------------------------------------------------
\section{Output Quality Hardening}

\subsection{Fix Warnings in Templates}

\begin{itemize}
\item Prefix unused \texttt{search\_pattern} with \texttt{\_} when entity
  has no search fields configured.
\item Remove unused imports from auth.rs (\texttt{ServiceRequest}, \texttt{HttpMessage}
  if not used).
\item Remove system-managed fields (shipped\_at, delivered\_at, last\_login\_at)
  from Create/Update payloads --- they should only be set by the system.
\item Add \texttt{\#[allow(dead\_code)]} to fields that exist for future use.
\end{itemize}

\subsection{Improve Service Templates}

Even without LLM enrichment, services should have:
\begin{itemize}
\item \texttt{require\_role(user, "level")?} at the start
\item \texttt{tracing::info!} logging with context
\item Proper \texttt{Option} handling for pagination params
\item \texttt{AppError::ValidationError(vec![msg.into()])} (correct constructor)
\end{itemize}

\subsection{Verification: All Three Domains Compile}

\begin{lstlisting}[language=bash]
cargo build --workspace --release && cargo test --workspace

for domain in warehouse ecommerce project; do
  rm -rf output/$domain
  cargo run --release -p isls-cli -- forge-v2 \
    --requirements examples/$domain.toml \
    --mock-oracle --passes 5 --output ./output/$domain
  cd output/$domain/backend && cargo build && cd ../../..
done

# Warnings per domain: target <= 5
\end{lstlisting}

% %%%%%%%%%%%%%%%%%%%%%%%%%%%%%%%%%%%%%%%%%%%%%%%%%%%%%%%%%%%%
\part{Self-Defining Norm System}
\label{part:norms}
% %%%%%%%%%%%%%%%%%%%%%%%%%%%%%%%%%%%%%%%%%%%%%%%%%%%%%%%%%%%%

\section{Norm Type System}

\subsection{New Crate: isls-norms}

Pure data crate. No IO, no network, no async. Contains all norm
types, the registry, composition engine, and built-in norms.

\begin{lstlisting}[language=C]
pub struct Norm {
    pub id: String,              // "ISLS-NORM-0042"
    pub name: String,            // "CRUD-Entity"
    pub level: NormLevel,
    pub triggers: Vec<TriggerPattern>,
    pub layers: NormLayers,
    pub parameters: Vec<NormParameter>,
    pub requires: Vec<String>,   // dependency norm IDs
    pub variants: Vec<NormVariant>,
    pub version: String,
    pub evidence: NormEvidence,
}

pub enum NormLevel {
    Atom,       // Field, Endpoint, Component
    Molecule,   // CRUD-Entity, JWT-Auth, Pagination
    Organism,   // Warehouse, ECommerce, ProjectTracker
    Ecosystem,  // SaaS, Marketplace
}

pub struct TriggerPattern {
    pub keywords: Vec<String>,
    pub concepts: Vec<String>,
    pub min_confidence: f64,
    pub excludes: Vec<String>,
}
\end{lstlisting}

\subsection{Norm Layers --- Cross-Layer Artifact Definitions}

\begin{lstlisting}[language=C]
pub struct NormLayers {
    pub database: Vec<DatabaseArtifact>,
    pub model: Vec<ModelArtifact>,
    pub query: Vec<QueryArtifact>,
    pub service: Vec<ServiceArtifact>,
    pub api: Vec<ApiArtifact>,
    pub frontend: Vec<FrontendArtifact>,
    pub test: Vec<TestArtifact>,
    pub config: Vec<ConfigArtifact>,
}

pub struct ModelArtifact {
    pub struct_name: String,
    pub fields: Vec<FieldSpec>,
    pub derives: Vec<String>,
    pub validations: Vec<ValidationSpec>,
}

pub struct FieldSpec {
    pub name: String,
    pub rust_type: String,
    pub sql_type: String,
    pub nullable: bool,
    pub default_value: Option<String>,
    pub indexed: bool,
    pub unique: bool,
    pub source: FieldSource,   // UserInput | SystemGenerated | SystemComputed
    pub description: String,
}

pub enum FieldSource { UserInput, UserOptional, SystemGenerated, SystemComputed }

pub struct ApiArtifact {
    pub method: String,
    pub path: String,
    pub auth_required: bool,
    pub min_role: String,
    pub request_body: Option<String>,
    pub response_type: String,
}

pub struct FrontendArtifact {
    pub component_type: FrontendComponent, // Page, Table, Form, Modal, Chart, etc.
    pub api_calls: Vec<String>,
}

pub enum FrontendComponent {
    Page, Table, Form, DetailView, Modal, DashboardCard,
    Chart, SearchBar, StatusBadge, ActionButton, Scanner,
}
\end{lstlisting}

\subsection{Norm Parameters and Variants}

\begin{lstlisting}[language=C]
pub struct NormParameter {
    pub name: String,
    pub param_type: ParamType,
    pub default: Option<String>,
}

pub enum ParamType {
    EntityName, FieldList, RelationshipList,
    RoleRequirement, Boolean, Choice(Vec<String>),
}

pub struct NormVariant {
    pub id: String,
    pub name: String,
    pub modifications: Vec<NormModification>,
}

pub enum NormModification {
    AddField(FieldSpec),
    AddEndpoint(ApiArtifact),
    AddComponent(FrontendArtifact),
    AddTest(TestArtifact),
    AddMigration(String),
    ModifyService(String, String),
}
\end{lstlisting}

% ----------------------------------------------------------
\section{Built-In Norm Catalog}

\subsection{Level 1 Molecules (20+)}

\begin{center}
\scriptsize
\begin{tabular}{llp{6.5cm}}
\toprule
\textbf{ID} & \textbf{Name} & \textbf{Layers} \\
\midrule
0042 & CRUD-Entity & Model+Create+Update+Validate+Migration+5Queries+5Service+5API+Table+Form+5Tests \\
0042-SD & CRUD+Soft-Delete & +is\_active field, filtered queries, restore endpoint \\
0042-AT & CRUD+Audit-Trail & +audit table, log trigger, audit endpoint \\
0088 & JWT-Auth & User+password hash+Claims+tokens+AuthUser extractor+login+register+roles \\
0100 & Pagination & Params+Response+Meta+LIMIT/OFFSET+page controls \\
0101 & Error-System & AppError enum+ResponseError+From conversions+JSON errors \\
0112 & Inventory & quantity+reorder\_level+StockMovement+adjust\_stock+reorder alert \\
0120 & State-Machine & status field+transitions map+update\_status+validation+badge \\
0130 & Search & ILIKE query+search input+debounced API call \\
0140 & Filter & WHERE per field+dropdown UI+query params \\
0150 & Sort & ORDER BY+column click+indicator \\
0156 & Barcode & unique sku/barcode+lookup endpoint+scanner component \\
0160 & File-Upload & file field+multipart endpoint+storage+preview \\
0170 & Export & CSV/JSON endpoint+download button \\
0180 & Dashboard-KPI & aggregation query+KPI card+auto-refresh \\
0190 & Chart & time-series query+canvas chart (bar/line/pie) \\
0200 & Notification & alert entity+condition trigger+list component+badge \\
0210 & Config & AppConfig struct+env loading+.env.example \\
0220 & Health-Check & /health+DB ping+uptime+version \\
0230 & Docker-Deploy & Dockerfile+compose+.env.example \\
0240 & README & Setup+API docs+architecture \\
\bottomrule
\end{tabular}
\end{center}

\subsection{Level 2 Organisms}

\begin{center}
\scriptsize
\begin{tabular}{llp{8cm}}
\toprule
\textbf{ID} & \textbf{Name} & \textbf{Composed Of} \\
\midrule
0500 & Warehouse & 0042$\times$7+0088+0112+0120+0156+0180$\times$4+0190$\times$3+0170+0230 \\
0501 & E-Commerce & 0042$\times$9+0088+0120+0130+0140+0180$\times$4+0190$\times$3+0230 \\
0502 & Project-Tracker & 0042$\times$8+0088+0120$\times$2+0130+0180$\times$3+0190$\times$2+0230 \\
\bottomrule
\end{tabular}
\end{center}

% ----------------------------------------------------------
\section{Norm Composition Engine}

\begin{lstlisting}[language=C]
/// Compose multiple activated norms into a generation plan.
pub fn compose_norms(
    norms: &[ActivatedNorm],
    params: &HashMap<String, String>,
) -> Result<ComposedPlan> {
    let mut plan = ComposedPlan::new();
    // 1. Instantiate each norm with parameters
    // 2. Merge layer artifacts (dedup by name)
    // 3. Resolve dependencies (activate required norms)
    // 4. Apply cross-norm wiring (order+inventory -> fulfill_order calls adjust_stock)
    // 5. Detect conflicts
    // 6. Generate interface contracts
    Ok(plan)
}

pub struct ActivatedNorm {
    pub norm: Norm,
    pub confidence: f64,
    pub source: ActivationSource, // UserExplicit, KeywordMatch, Dependency, ChatAmendment
}

pub struct ComposedPlan {
    pub database: Vec<DatabaseArtifact>,
    pub models: Vec<ModelArtifact>,
    pub queries: Vec<QueryArtifact>,
    pub services: Vec<ServiceArtifact>,
    pub api: Vec<ApiArtifact>,
    pub frontend: Vec<FrontendArtifact>,
    pub tests: Vec<TestArtifact>,
    pub config: Vec<ConfigArtifact>,
    pub interfaces: Vec<InterfaceContract>,
}
\end{lstlisting}

\subsection{Cross-Norm Wiring}

\begin{lstlisting}[language=C]
pub struct NormWiring {
    pub when: (String, String), // (norm_a_id, norm_b_id)
    pub add_services: Vec<ServiceArtifact>,
    pub add_rules: Vec<BusinessRule>,
    pub add_tests: Vec<TestArtifact>,
}

// Example: when Order + Inventory both active,
// add fulfill_order that calls adjust_stock
fn wiring_order_inventory() -> NormWiring { ... }
// Example: when Auth active, all API endpoints get auth check
fn wiring_auth_all_endpoints() -> NormWiring { ... }
\end{lstlisting}

\subsection{NormRegistry}

\begin{lstlisting}[language=C]
pub struct NormRegistry {
    norms: HashMap<String, Norm>,
    wirings: Vec<NormWiring>,
    candidates: HashMap<String, NormCandidate>,
    auto_id_counter: u32,
    persistence_path: Option<PathBuf>,
}

impl NormRegistry {
    pub fn default() -> Self; // load built-in norms
    pub fn match_description(&self, desc: &str) -> Vec<ActivatedNorm>;
    pub fn get(&self, id: &str) -> Option<&Norm>;
    pub fn wirings_for(&self, a: &str, b: &str) -> Vec<&NormWiring>;
    pub fn save(&self) -> Result<()>;
    pub fn load(&mut self) -> Result<()>;
}
\end{lstlisting}

% ----------------------------------------------------------
\section{Self-Defining: Crystal $\to$ Norm Promotion}

\subsection{Cross-Layer Pattern Detection}

After each forge run, Barbara analyzes the output:

\begin{lstlisting}[language=C]
pub struct CrossLayerPattern {
    pub signature: String,          // SHA-256
    pub layers_present: Vec<LayerType>,
    pub artifacts: Vec<ObservedArtifact>,
    pub observed_on: String,        // entity name
    pub domain: String,
    pub run_id: String,
}

pub struct ObservedArtifact {
    pub layer: LayerType,
    pub artifact_type: String,
    pub name: String,
    pub signature: String,
    pub field_names: Vec<String>,
}
\end{lstlisting}

\subsection{NormCandidate}

\begin{lstlisting}[language=C]
pub struct NormCandidate {
    pub id: String,                 // ISLS-CAND-XXXX
    pub observations: Vec<CrossLayerPattern>,
    pub domains: Vec<String>,
    pub observation_count: usize,
    pub consistency: f64,           // avg topology similarity
    pub consistent_layers: Vec<LayerType>,
    pub common_artifacts: Vec<AbstractedArtifact>,
    pub status: CandidateStatus,    // Observing | Eligible | Promoted | Rejected
}

pub struct AbstractedArtifact {
    pub layer: LayerType,
    pub name_template: String,      // "{entity}_status"
    pub common_fields: Vec<AbstractedField>,
    pub confidence: f64,            // fraction of observations with this artifact
}
\end{lstlisting}

\subsection{Promotion Criteria}

\begin{lstlisting}[language=C]
pub struct PromotionCriteria {
    pub min_consistency: f64,       // 0.85
    pub min_domains: usize,         // 3
    pub min_layers: usize,          // 4
    pub min_observations: usize,    // 5
    pub min_artifact_presence: f64, // 0.80
}
\end{lstlisting}

A candidate meeting ALL thresholds is auto-promoted to a Norm.

\subsection{Automatic Norm Synthesis}

\begin{lstlisting}[language=C]
pub fn synthesize_norm(candidate: &NormCandidate) -> Norm {
    let id = format!("ISLS-NORM-AUTO-{:04}", next_auto_id());
    let name = infer_norm_name(&candidate.common_artifacts);
    let triggers = infer_triggers(&candidate.observations);
    let layers = build_norm_layers(&candidate.common_artifacts);
    let parameters = infer_parameters(&candidate.observations);
    Norm { id, name, level: NormLevel::Molecule, triggers, layers, parameters, ... }
}
\end{lstlisting}

\subsection{The Learning Loop}

\begin{lstlisting}[language=C]
impl NormRegistry {
    pub fn observe_and_learn(
        &mut self,
        artifacts: &[Artifact],
        reader: &Reader,
        topo: &CodeTopoAnalyzer,
        domain: &str,
        run_id: &str,
    ) {
        let patterns = extract_cross_layer_patterns(artifacts, reader, topo);
        for pattern in &patterns {
            match self.match_pattern_to_norm(pattern) {
                Some(norm) => self.reinforce_norm(&norm.id, pattern, domain),
                None => self.add_to_candidate_pool(pattern, domain, run_id),
            }
        }
        let promoted = self.check_promotions();
        for id in &promoted {
            let norm = synthesize_norm(self.candidates.get(id).unwrap());
            self.register(norm);
        }
    }
}
\end{lstlisting}

\subsection{Self-Bootstrap Trajectory}

\begin{lstlisting}
Run 0:  20 norms. Generate warehouse. 4 novel candidates.
Run 1:  20 norms, 4 candidates. Generate ecommerce. 3 match, 2 new.
Run 2:  20 norms, 6 candidates. Generate project tracker.
        Candidate "State-Machine": 3 domains, 5 obs, 0.91 -> PROMOTED.
Run 3:  21 norms (20+1 auto). Booking system. Auto-norm activates.
Run 10: 24 norms (20+4 auto).
Run 50: 35 norms (20+15 auto). System knows patterns no human defined.
\end{lstlisting}

\subsection{Persistence}

Auto-norms and candidates saved to \texttt{\textasciitilde/.isls/norms.json}.
Built-in norms are always in memory. Deleting the file resets
auto-discovery but retains the 20 built-in norms.

% %%%%%%%%%%%%%%%%%%%%%%%%%%%%%%%%%%%%%%%%%%%%%%%%%%%%%%%%%%%%
\part{Chat Agent}
\label{part:chat}
% %%%%%%%%%%%%%%%%%%%%%%%%%%%%%%%%%%%%%%%%%%%%%%%%%%%%%%%%%%%%

\section{New Crate: isls-chat}

\subsection{Intent Recognition}

\begin{lstlisting}[language=C]
pub struct ChatIntent {
    pub intent_type: IntentType,
    pub entities: Vec<String>,
    pub fields: Vec<FieldIntent>,
    pub features: Vec<String>,
    pub rules: Vec<RuleIntent>,
    pub roles: Vec<String>,
    pub raw_text: String,
}

pub enum IntentType {
    CreateApplication,  // "I need an app that..."
    AddField,           // "Add weight to parts"
    RemoveFeature,      // "Remove barcode scanning"
    ModifyEntity,       // "Change order statuses to..."
    AddBusinessRule,    // "Warn when stock is low"
    ViewApplication,    // "Show me the project"
    Deploy,             // "Deploy it"
    Help,               // "What can this do?"
    Clarify,            // Unclear
}
\end{lstlisting}

\subsection{Intent Extraction via LLM}

\begin{lstlisting}[language=C]
pub fn extract_intent(oracle: &dyn Oracle, message: &str) -> Result<ChatIntent> {
    let prompt = format!(
        "Extract software concepts from this user message.\n\
         Return JSON only, no markdown:\n\
         {{\"intent\": \"create|add_field|modify|add_rule|help\",\n\
           \"entities\": [\"list of data types\"],\n\
           \"features\": [\"list of features\"],\n\
           \"fields\": [{{\"entity\": \"X\", \"name\": \"Y\", \"type\": \"Z\"}}],\n\
           \"rules\": [{{\"trigger\": \"X\", \"condition\": \"Y\", \"action\": \"Z\"}}],\n\
           \"roles\": [\"list of user roles\"]}}\n\n\
         Message: {}", message);
    let json = oracle.call_json(&prompt, 1024)?;
    parse_intent(json)
}

/// Fallback for --mock-oracle: keyword-based intent recognition.
pub fn extract_intent_keywords(message: &str) -> ChatIntent {
    // "I need an app" -> CreateApplication
    // "add" + field name -> AddField
    // "remove" / "delete" -> RemoveFeature
    // Simple keyword matching, no LLM needed
}
\end{lstlisting}

\subsection{Intent to Norm Mapping}

\begin{lstlisting}[language=C]
pub fn intent_to_norm_ops(
    intent: &ChatIntent,
    registry: &NormRegistry,
) -> Vec<NormOperation> {
    match intent.intent_type {
        IntentType::CreateApplication => {
            let norms = registry.match_description(&intent.raw_text);
            vec![NormOperation::ComposeNew(norms)]
        }
        IntentType::AddField => {
            intent.fields.iter().map(|f| NormOperation::AmendEntity {
                entity: f.entity.clone(),
                add_field: field_from_intent(f),
            }).collect()
        }
        IntentType::AddBusinessRule => {
            intent.rules.iter().map(|r| NormOperation::AddRule {
                trigger: r.trigger.clone(),
                condition: r.condition.clone(),
                action: r.action.clone(),
            }).collect()
        }
        _ => vec![],
    }
}

pub enum NormOperation {
    ComposeNew(Vec<ActivatedNorm>),
    AmendEntity { entity: String, add_field: FieldSpec },
    AddNorm(String),
    RemoveNorm(String),
    AddRule { trigger: String, condition: String, action: String },
}
\end{lstlisting}

\subsection{Incremental Regeneration}

\begin{lstlisting}[language=C]
/// After a plan amendment, regenerate ONLY affected files.
pub fn affected_files(op: &NormOperation) -> Vec<String> {
    match op {
        NormOperation::AmendEntity { entity, .. } => vec![
            format!("backend/src/models/{}.rs", entity.to_snake()),
            format!("backend/src/database/{}_queries.rs", entity.to_snake()),
            format!("backend/src/services/{}.rs", entity.to_snake()),
            format!("backend/src/api/{}.rs", entity.to_snake()),
            "backend/database/migrations/001_initial.sql".into(),
            format!("frontend/src/pages/{}.js", entity.to_snake()),
        ],
        _ => vec![],
    }
}
\end{lstlisting}

% %%%%%%%%%%%%%%%%%%%%%%%%%%%%%%%%%%%%%%%%%%%%%%%%%%%%%%%%%%%%
\part{Studio Dashboard}
\label{part:studio}
% %%%%%%%%%%%%%%%%%%%%%%%%%%%%%%%%%%%%%%%%%%%%%%%%%%%%%%%%%%%%

\section{Architecture}

Single HTML file. No npm. No React. Vanilla JS + fetch + CSS grid.
Served by \texttt{isls-gateway} at \texttt{http://localhost:3000}.

\subsection{Layout}

\begin{lstlisting}
+-------+-----------------------------------+-------------+
| CHAT  |          MAIN PANEL               |   NORM      |
|       |  [Projects] [Runs] [Files]        |  EXPLORER   |
| User  |                                   |             |
| types |  Project: warehouse-system        |  Active:    |
| here  |  Status: Generated                |  0042 x7    |
|       |  Files: 60 | LOC: 5120            |  0088       |
| ISLS  |  Crystals: 14 | Evidence: OK      |  0112       |
| resp- |                                   |  ...        |
| onds  |  [Convergence Graph]              |             |
|       |  [Crystal Stats Table]            |  Candidates:|
| [...] |                                   |  CAND-005   |
| Input |                                   |  [+Norm]    |
+-------+-----------------------------------+-------------+
\end{lstlisting}

\subsection{API Endpoints (isls-gateway extensions)}

\begin{lstlisting}[language=C]
// Chat
POST   /api/chat              { message: String } -> ChatResponse
GET    /api/chat/history       -> Vec<ChatMessage>

// Projects
GET    /api/projects           -> Vec<ProjectSummary>
POST   /api/projects/forge     { description: String, settings: Settings } -> ForgeResult
GET    /api/projects/:id       -> ProjectDetail
DELETE /api/projects/:id

// Runs
GET    /api/projects/:id/runs  -> Vec<RunSummary>
GET    /api/runs/:id/trace     -> DecompositionTrace

// Files
GET    /api/projects/:id/files          -> Vec<FileInfo>
GET    /api/projects/:id/files/*path    -> FileContent

// Norms
GET    /api/norms              -> Vec<NormSummary>
GET    /api/norms/:id          -> NormDetail
GET    /api/norms/candidates   -> Vec<CandidateSummary>

// Crystals
GET    /api/crystals           -> Vec<CrystalSummary>
GET    /api/crystals/stats     -> CrystalStats

// System
GET    /api/health             -> HealthStatus

pub struct ChatResponse {
    pub message: String,
    pub intent: String,
    pub norms_activated: Vec<String>,
    pub estimated_files: usize,
    pub estimated_loc: usize,
    pub action_buttons: Vec<ActionButton>,
}

pub struct ActionButton {
    pub label: String,
    pub action: String,
    pub style: String,
}
\end{lstlisting}

\subsection{Chat Flow}

\begin{lstlisting}
1. User types -> POST /api/chat
2. System returns: explanation + norms + estimation + buttons
3. User clicks [Generate] -> POST /api/projects/forge
4. Progress: Pass 0 done, Pass 1 running, ...
5. Done: project in dashboard, files browsable
6. User chats: "Add weight to parts" -> incremental regen
\end{lstlisting}

% %%%%%%%%%%%%%%%%%%%%%%%%%%%%%%%%%%%%%%%%%%%%%%%%%%%%%%%%%%%%
\part{Implementation Plan}
\label{part:impl}
% %%%%%%%%%%%%%%%%%%%%%%%%%%%%%%%%%%%%%%%%%%%%%%%%%%%%%%%%%%%%

\section{Execution Order}

\begin{enumerate}
\item \textbf{Part I --- TypeContext + Enrichment} (isls-renderloop).
  Test: LLM enrichment compiles with correct types.
\item \textbf{Part I --- Multi-Domain} (isls-hypercube + orchestrator).
  Test: all 3 domains produce compilable backends.
\item \textbf{Part I --- Quality Hardening} (orchestrator templates).
  Test: $\leq$ 5 warnings per domain.
\item \textbf{Part II --- isls-norms} types + registry + 20 molecules + 3 organisms + composition.
  Test: compose warehouse norms, verify 7 entities.
\item \textbf{Part II --- Self-defining} candidate pool + promotion + synthesis.
  Test: 3 mock runs, pattern promoted, activates in 4th.
\item \textbf{Part II --- Wire norms into pipeline.}
  TOML $\to$ norm matching $\to$ composition $\to$ decomposer.
  Test: forge-v2 still works, now via norm path.
\item \textbf{Part III --- isls-chat} intent recognition + mapping + incremental regen.
  Test: ``scan parts'' extracts 4 entities + barcode feature.
\item \textbf{Part IV --- isls-gateway} API endpoints.
  Test: POST /api/chat returns norms.
\item \textbf{Part IV --- isls-studio} single HTML dashboard.
  Test: browser shows chat + dashboard + norm explorer.
\item \textbf{Part IV --- isls-agent} wire chat+norms+forge.
  Test: end-to-end chat $\to$ compiled output.
\item \textbf{Part IV --- Incremental regen} via chat amendment.
  Test: add field, only 6 files change, still compiles.
\end{enumerate}

Commit after each numbered step.

\section{Success Criteria}

\begin{enumerate}
\item \texttt{cargo build --workspace} --- zero errors at every step.
\item All existing tests pass at every step.
\item 3 TOML examples produce compilable backends in mock mode.
\item Warning count $\leq$ 5 per generated backend.
\item With \texttt{--api-key}: LLM-enriched services compile
  (hallucinated fields rejected by TypeContext validation).
\item NormRegistry has $\geq$ 20 molecules + 3 organisms.
\item Self-discovery: after 3 mock runs, $\geq$ 1 candidate exists.
  After criteria met, candidate promoted to auto-norm.
\item \texttt{forge-v2 --requirements X.toml} works via norm path
  (backward compatible).
\item POST \texttt{/api/chat} returns correct norms + estimation.
\item Studio at \texttt{http://localhost:3000} shows chat + dashboard.
\item Chat-to-generation flow works end-to-end.
\item Chat amendment regenerates only affected files.
\item Crystal stats visible in Studio.
\item Norm explorer shows built-in (blue), auto (green), candidates (yellow).
\end{enumerate}

\section{Constraints for Claude Code}

\begin{enumerate}
\item \textbf{Repository: graphity/isls.} All work in the isls workspace.
\item \textbf{v2.1 must not break.} \texttt{forge-v2 --mock-oracle}
  without new flags produces same output.
\item \textbf{isls-norms is pure data.} No IO, no network. Compiles
  in $<$ 1s.
\item \textbf{Studio is ONE HTML file.} No npm. No React. Vanilla JS.
  Served via \texttt{include\_str!} or static dir.
\item \textbf{Chat uses Oracle for intent extraction.} With
  \texttt{--mock-oracle}, use keyword matching.
\item \textbf{Norms are Rust structs in the binary.} Not JSON/TOML
  files. Auto-norms persist in \texttt{\textasciitilde/.isls/norms.json}.
\item \textbf{Sequential.} Part I complete before Part II. Part II
  before Part III. Part III before Part IV.
\item \textbf{Commit after each step.} 11 commits minimum.
\item \textbf{The chat-to-generation flow is the most important
  deliverable.} If time runs out: skip incremental regen, skip
  norm explorer. But chat $\to$ generate must work.
\item \textbf{No \texttt{unwrap()} in library code.}
\item \textbf{Every public item documented.}
\item \textbf{MIT License. Copyright (c) 2026 Sebastian Klemm.}
\end{enumerate}

\vfill
\noindent\rule{\textwidth}{0.4pt}
\begin{center}
\small
ISLS v2.2--v3.0 Unified Specification --- Sebastian Klemm --- March 2026\\
Type-aware enrichment. Three domains. Self-defining norms.\\
Chat-driven generation. Studio dashboard.\\
Homer types. The system builds. The norms grow.\\
\url{https://github.com/LashSesh/graphity}
\end{center}

\end{document}
